\section{Záver}

V práci bol navrhnutý a implementovaný systém na zber časov stlačenia kláves
a na identifikáciu používateľa z jedného pokusu. Webová aplikácia zaznamenáva
stlačenia a uvoľnenia kláves, serverová časť ich ukladá a z prvých desiatich
pokusov každého účastníka vytvára štatistickú šablónu. Nový pokus sa porovná
so šablónami pomocou štandardizovanej euklidovskej vzdialenosti. Okrem dlhej vety
systém podporuje aj režim s jedným krátkym slovom a porovnanie oboch režimov.
Klientska aplikácia bola overená funkčným testom a systém je dostupný
v repozitári \cite{repo}.

Na dátach piatich účastníkov dosiahol systém na 33 testovacích pokusoch
úspešnosť top-1 78,79\,\% a top-3 96,97\,\%. Vzorka je však malá a nerovnomerná,
preto treba výsledok chápať ako orientačný.

% ----- Zhodnotenie splnenia cieľa -----
Cieľ práce sa podarilo splniť. Vznikol funkčný systém, ktorý zaznamená rytmus
písania, uloží ho, vytvorí šablóny účastníkov a z jedného nového pokusu určí
najpodobnejšieho účastníka alebo ho označí ako neznámeho. Úspešnosť top-1
78,79\,\% je výrazne vyššia ako náhodný odhad pri piatich účastníkoch (20\,\%)
a úspešnosť top-3 96,97\,\% je vyššia ako náhodných 60\,\%. K riešeniu patrí aj
príručka na jeho používanie a zdrojový kód v repozitári, takže postup možno
zopakovať.

% ----- Obmedzenia riešenia -----
Výsledky treba brať ako orientačné. Vzorka je malá (päť účastníkov a 33
testovacích pokusov) a nerovnomerná: jeden účastník prispel 15 pokusmi, iný len
tromi. Všetky údaje vznikli na notebooku, preto sa nedá posúdiť vplyv zariadenia
a vplyv klávesnice len veľmi obmedzene (dve relácie s externou USB
klávesnicou). Režim slova má zatiaľ dáta len od jedného účastníka, preto
porovnanie vety a slova ešte nie je vyhodnotené. Prah pre označenie neznámeho pokusu je
zvolený empiricky a nebol kalibrovaný na samostatných dátach. Šablóny sa
počítajú pri každej požiadavke nanovo, čo stačí pre malý súbor dát, no pre
väčší by to bolo pomalé. Identifikátorom účastníka je meno, takže databáza
obsahuje osobné údaje.

% ----- Smery ďalšieho vývoja -----
V ďalšom vývoji je vhodné zozbierať dáta od viacerých účastníkov v niekoľkých
reláciách, na rôznych zariadeniach a klávesniciach, aby sa dal vyhodnotiť ich
vplyv. Po zbere dát od viacerých účastníkov v oboch režimoch treba vyhodnotiť
porovnanie vety a slova. Prah
neznámeho pokusu by sa mal kalibrovať na samostatnej validačnej množine a
šablóny by sa mali ukladať namiesto opakovaného výpočtu. Namiesto mena by sa
mal používať pseudonymný kód účastníka. Zaujímavé by bolo porovnať zvolenú
vzdialenosť s inými mierami podobnosti a vyhodnotiť aj scenár overovania
identity, v ktorom sa sleduje prijatie cudzieho a odmietnutie vlastného
používateľa.
